\section{Conclusion}
\label{sec:conclusion}

For the systems tested using \ac{fmri} responses to English-language stimuli, neither intervention route established a participant-general training benefit attributable to brain-response content.\evd{E008}\evd{E025}\evd{E026}\evd{E030}
The recorded-response route supported controlled brain predictivity, but did not demonstrate, across the tested objectives and variants, that correctly paired responses beat permuted responses for participants in general.\evd{E002}\evd{E006}\evd{E008}\evd{E011}\evd{E013}\evd{E017}
The WikiText checks in the synthetic-response route supported target uptake (the target became easier to recover from the student) and retained-student movement (the retained student still differed from its baseline), both at acceptable measured language quality.\evd{E016}\evd{E025}
However, the target-projection assay did not meet its continuation rule, and the text-derived auxiliary-control arm was non-identifying, too poorly matched to attribute any difference to brain-response content.\evd{E026}\evd{E030}
Direct participant-level biological transfer over ordinary \ac{kd} was also not demonstrated.\evd{E025}
The utility endpoint for recorded-response training remained unresolved, and brain-specific utility was not demonstrated.\evd{E009}
The tested routes therefore do not establish brain alignment as a usable training signal, but they do not rule out a smaller benefit in these systems or a benefit from other targets, interventions, or populations.
The supported contribution is narrower than the governing question: controlled brain predictivity in the declared assays, a demonstrated ability to change the student through synthetic-response training, and an evidence framework that keeps those findings distinct from attribution, biological transfer, and utility.

A controlled brain predictivity score alone does not establish a usable training signal.
Within the evidence criteria used in this thesis, target uptake and retained-student movement must pass before transfer or attribution is evaluated, and a utility claim framed as the payoff of improved controlled brain predictivity requires that prerequisite first.
After optimization, the following outcomes must be evaluated separately: target uptake and retained-student movement at acceptable measured language quality; direct participant-level transfer over ordinary \ac{kd}; brain-response-specific attribution against an adequate route-appropriate control; and any claimed utility endpoint.
For recorded-response training, a permuted-response control tests whether correct stimulus--response pairing matters.
For synthetic-response training, an adequately matched text-derived auxiliary-control arm tests whether any observed benefit is specific to brain-response content.
Seeds, folds, voxels, and \acp{roi} cannot substitute for participants when the claim concerns people.
Keeping these claim dependencies separate shows what an intervention achieved and what evidence remains necessary before controlled brain predictivity can be treated as a usable training~signal.
